% =============================================================
%  CS102-F26 Digital Design and Computer Architecture I
%  Lecture 20 -- Processor Design (Unit 20)
%  Follows the introduction of Section 6.6 and Section 6.6.1 of
%  A Student's Computer Architecture Book I.
%
%  Build the slides and the print handout:  npm run materials -- lecture-20
% =============================================================
\makeatletter
\def\input@path{{templates/latex/}{../}}
\makeatother

\documentclass[10pt]{wosm-slides}
\usepackage{cs102-lectures}[2026/10/20]

\title[Lecture 20 \textperiodcentered\ Processor Design]{Processor Design}
\subtitle{Lecture 20: fetch, decode, execute and load; single-cycle, multi-cycle and
  pipelined designs}

% ---------- Helpers for this lecture --------------------------
\tikzset{
  stage/.style={draw=wosm-text, line width=0.5pt, fill=wosm-surface},
}
% One instruction on a time line, in clock cycles:
% \fdel{<y>}{<start>}{<fetch>}{<decode>}{<execute>}{<load>}{<label>}
\newcommand*\fdel[7]{%
  \pgfmathsetmacro\xa{#2}\pgfmathsetmacro\xb{\xa+#3}%
  \pgfmathsetmacro\xc{\xb+#4}\pgfmathsetmacro\xd{\xc+#5}%
  \pgfmathsetmacro\xe{\xd+#6}%
  \foreach \p/\q/\t in {\xa/\xb/F, \xb/\xc/D, \xc/\xd/E, \xd/\xe/L} {
    \draw[stage] (\p,#1+0.2) rectangle (\q,#1-0.2);
    \node[font=\tiny] at ({(\p+\q)/2},#1) {\t};
  }
  \node[left, font=\scriptsize] at (-0.15,#1) {#7};}
% A time axis in clock cycles: \cycles{<y>}{<last cycle>}{<step of the labels>}
\newcommand*\cycles[3]{%
  \draw[wosm-muted, line width=0.4pt] (0,#1) -- (#2,#1);
  \foreach \c in {0,#3,...,#2} {
    \draw[wosm-muted, line width=0.4pt] (\c,#1) -- (\c,#1-0.08);
    \node[below, font=\tiny, text=wosm-muted] at (\c,#1-0.05) {\c};
  }}

% =============================================================
\begin{document}

\maketitle

\begin{frame}{Today}
  \tableofcontents
\end{frame}

% =============================================================
\section{The four stages}

\begin{frame}{Building the DX8}
  \begin{columns}[c, totalwidth=\textwidth]
    \column{0.5\textwidth}
      The parts we have: the RAM and the register file \unitref{19}, the ALU
      \unitref{18}, and registers such as the IP, IR, regA, regB, regOUT and regFLAG
      \unitref{15}.\par\medskip
      Two more parts make it run:
      \begin{itemize}
        \item the \alert{stepper}, which turns on one output at a time, in a fixed
          order, in step with the clock;
        \item the \alert{controller}, which sets the control bits of every other part
          for the current step.
      \end{itemize}
    \column{0.44\textwidth}
      \begin{block}{The question for today}
        Running one instruction takes several things, one after another: read it,
        work out what it means, do it, and keep the result. How do we organise that
        in time?
      \end{block}
  \end{columns}
\end{frame}

\begin{frame}{Fetch, decode, execute, load}
  \begin{center}
    \small
    \begin{tabular}{@{}lp{0.76\textwidth}@{}}
      \toprule
      Stage & What happens \\
      \midrule
      \textbf{Fetch} & The controller sets RE of the RAM, and WE of the register file
        and the IR: the instruction is read from memory into the CPU. \\
      \textbf{Decode} & The instruction is examined: is it an ALU or a memory
        instruction, and which part is responsible for it? \\
      \textbf{Execute} & The instruction is carried out, in the ALU or elsewhere; the
        results are kept for now in registers such as regOUT and regFLAG. \\
      \textbf{Load} & The kept results are written back where they belong, to the
        register file or the RAM, if needed. Then the cycle starts again. \\
      \bottomrule
    \end{tabular}
  \end{center}

  \begin{remark}
    Splitting the work into four stages makes the design easier and modular: we can
    design and test one stage at a time.
  \end{remark}
\end{frame}

\begin{frame}{One instruction, four stages}
  \begin{columns}[T, totalwidth=\textwidth]
    \column{0.48\textwidth}
      \wosmlabel{\texttt{add}: regOUT $\leftarrow$ regA $+$ regB}
      \begin{enumerate}
        \item \textbf{Fetch}: the instruction at the address in the IP comes in
          from the RAM.
        \item \textbf{Decode}: $b_{15} = 1$, so the ALU is responsible; the opcode
          0001 goes to the ALU.
        \item \textbf{Execute}: the ALU adds; the sum goes to regOUT, the flags to
          regFLAG.
        \item \textbf{Load}: regOUT and regFLAG go back to the register file.
      \end{enumerate}
    \column{0.48\textwidth}
      \wosmlabel{\texttt{store}: $M[\texttt{memAddr}] \leftarrow$ regX}
      \begin{enumerate}
        \item \textbf{Fetch}: the same as for every instruction.
        \item \textbf{Decode}: $b_{15} = 0$ and opcode 0001: a \texttt{store}.
        \item \textbf{Execute}: there is nothing for the ALU to do.
        \item \textbf{Load}: the register's content is written into the RAM; the
          book's DX8 gives the RAM its address through the IP.
      \end{enumerate}
  \end{columns}
\end{frame}

% =============================================================
\section{Three ways to run the stages}

\begin{frame}{Single-cycle and multi-cycle}
  \begin{center}
    \begin{tikzpicture}[x=0.4cm, y=0.82cm]
      \node[anchor=west, font=\scriptsize] at (0,1.0) {\textbf{Single-cycle}: one
        long clock cycle per instruction};
      \foreach \i in {1,2,3} {
        \fdel{0.3}{8*(\i-1)}{2}{2}{2}{2}{}
        \node[font=\tiny, text=wosm-muted] at (8*\i-4,0.7) {instr.~\i};
      }
      \foreach \c in {0,8,16,24} {\draw[wosm-brick, line width=0.6pt] (\c,0.6) --
        (\c,0.0);}
      \cycles{-0.05}{24}{4}
      \node[anchor=west, font=\scriptsize] at (0,-0.85) {\textbf{Multi-cycle}: short
        cycles, as many as each stage needs};
      \fdel{-1.5}{0}{2}{2}{2}{2}{instr.~1}
      \fdel{-2.1}{8}{1}{1}{2}{1}{instr.~2}
      \fdel{-2.7}{13}{2}{2}{1}{2}{instr.~3}
      \cycles{-3.1}{24}{4}
    \end{tikzpicture}
  \end{center}

  \begin{itemize}
    \item Single-cycle: the period must fit the slowest instruction. If that takes
      100\,ns, \emph{every} instruction takes 100\,ns, even one that could finish
      sooner.
    \item Multi-cycle: here the three instructions take 8, 5 and 7 short cycles: 20
      instead of 24. But each instruction still waits for the one before it.
  \end{itemize}
\end{frame}

\begin{frame}{Pipelined}
  \begin{center}
    \begin{tikzpicture}[x=0.4cm]
      \node[anchor=west, font=\scriptsize] at (0,0.7) {\textbf{Multi-cycle}, 2 cycles
        per stage: 24 cycles for 3 instructions};
      \foreach \i in {1,2,3} {
        \fdel{0.05}{8*(\i-1)}{2}{2}{2}{2}{}
        \node[font=\tiny, text=wosm-muted] at (8*\i-4,0.42) {instr.~\i};
      }
      \cycles{-0.25}{24}{2}
      \node[anchor=west, font=\scriptsize] at (0,-1.05) {\textbf{Pipelined}: 12 cycles
        for the same 3};
      \fdel{-1.6}{0}{2}{2}{2}{2}{instr.~1}
      \fdel{-2.1}{2}{2}{2}{2}{2}{instr.~2}
      \fdel{-2.6}{4}{2}{2}{2}{2}{instr.~3}
      \cycles{-3.0}{24}{2}
    \end{tikzpicture}
  \end{center}

  \begin{itemize}
    \item When instruction 1 moves on to decode, the fetch stage is idle. A pipelined
      CPU fetches instruction 2 at once: the stages work on different instructions
      at the same time.
    \item Each instruction still takes 8 cycles; but one finishes every 2 cycles.
  \end{itemize}
\end{frame}

\begin{frame}{Latency and throughput}
  \begin{columns}[T, totalwidth=\textwidth]
    \column{0.5\textwidth}
      \begin{definition}[Latency and throughput]
        \emph{Latency}: the time a CPU takes to finish one instruction.
        \emph{Throughput}: how many instructions it finishes per unit of time.
      \end{definition}
      \begin{itemize}
        \item Single-cycle to multi-cycle: shorter latency, for instructions that
          need fewer cycles.
        \item Multi-cycle to pipelined: the same latency, more throughput.
      \end{itemize}
    \column{0.46\textwidth}
      \centering
      \small
      \begin{tabular}{@{}lcc@{}}
        \toprule
        & Multi-cycle & Pipelined \\
        \midrule
        Latency & 8 & 8 \\
        3 instructions & 24 & 12 \\
        Per cycle & $\frac{3}{24} = 0.125$ & $\frac{3}{12} = 0.25$ \\
        $n$ instructions & $8n$ & $2(n + 3)$ \\
        \bottomrule
      \end{tabular}\par\medskip
      \wosmlabel{4 stages of 2 cycles each. For long programs the pipeline finishes
      one instruction every 2 cycles: 4 times the throughput.}
  \end{columns}
\end{frame}

\begin{frame}{The catch}
  \begin{itemize}
    \item A pipeline needs registers between its stages, so that each stage can hold
      its own instruction: the ``register in the middle of a long path'' of
      \unitref{14}.
    \item An instruction may need a result that the instruction just before it has
      not written back yet.
    \item After a jump, the instructions already fetched behind it may be the wrong
      ones.
    \item Handling these cases is what makes pipelined CPUs the hardest of the three
      to design.
  \end{itemize}

  \begin{remark}
    These problems are called \emph{hazards}. They are beyond this module: the DX8 is
    not pipelined, so it never meets them.
  \end{remark}
\end{frame}

\begin{frame}{The DX8's choice}
  \begin{center}
    \begin{tabular}{@{}lll@{}}
      \toprule
      Design & To build & Performance \\
      \midrule
      Single-cycle & simplest & lowest \\
      Multi-cycle & moderate & moderate \\
      Pipelined & hardest & highest \\
      \bottomrule
    \end{tabular}
  \end{center}

  \begin{itemize}
    \item The book chooses \alert{multi-cycle} for the DX8: a trade-off between
      performance and simplicity.
    \item Each stage gets at most 2 clock cycles, or \emph{steps}: fetch is steps 1
      and 2, decode 3 and 4, execute 5 and 6, load 7 and 8. Then back to step 1.
    \item The stepper counts the steps \unitref{21}.
  \end{itemize}
\end{frame}

\begin{frame}{Try it}
  \begin{enumerate}
    \item In a single-cycle CPU, the slowest instruction takes 100\,ns. What is the
      fastest clock, and how long does the quickest instruction take?\par
      \uncover<2->{\small 10\,MHz; 100\,ns, like every other instruction.}
    \item The DX8 takes 8 cycles per instruction. With a 1\,MHz clock, how many
      instructions does it finish per second?\par
      \uncover<3->{\small $10^6 / 8 = 125\,000$.}
    \item 4 stages of 2 cycles each. How many cycles do 10 instructions take,
      multi-cycle and pipelined?\par
      \uncover<4->{\small $8 \cdot 10 = 80$, and $2 \cdot (10 + 3) = 26$.}
    \item Does pipelining make one instruction finish sooner?\par
      \uncover<5->{\small No: its latency stays 8 cycles. More instructions finish per
        second.}
  \end{enumerate}
\end{frame}

% =============================================================
\section{Summary}

\begin{frame}{Summary}
  \begin{itemize}
    \item Every instruction goes through four stages: fetch, decode, execute and
      load.
    \item Single-cycle: all four in one long clock cycle, set by the slowest
      instruction.
    \item Multi-cycle: shorter cycles; each stage takes the cycles it needs, one
      instruction at a time.
    \item Pipelined: the stages overlap, working on different instructions; the
      same latency, more throughput.
    \item The DX8 is multi-cycle: 8 steps per instruction, 2 for each stage.
  \end{itemize}
\end{frame}

\begin{frame}{Before the next lecture}
  \begin{itemize}
    \item Read the beginning of Section~6.6 and Section~6.6.1 of the book.
    \item Draw the time line of 5 instructions on the multi-cycle DX8 and on a
      pipelined DX8, and count the cycles.
    \item Read Sections~6.6.2 and 6.6.3.
  \end{itemize}

  \begin{block}{Next lecture \unitref{21}}
    The stepper, the data path and the controller.
  \end{block}
\end{frame}

\end{document}
